\documentclass[11pt, a4paper]{article}
\usepackage[a4paper, top=2.5cm, bottom=2.5cm, left=2cm, right=2cm]{geometry}
\usepackage{fontspec}
\usepackage[english, bidi=basic, provide=*]{babel}
\babelprovide[import, onchar=ids fonts]{english}
%\babelfont{rm}{Noto Sans}
\babelfont{rm}{TeX Gyre Termes}
\babelfont{sf}{TeX Gyre Heros}
\babelfont{tt}{TeX Gyre Cursor}


\usepackage{enumitem}
\usepackage{amsmath}

\title{\textbf{Chapter 1: Evolution, the Themes of Biology, and Scientific Inquiry \\ 50 Multiple Choice Questions}}
\author{}
\date{}

\begin{document}

\maketitle

\begin{enumerate}

\item \textbf{Which of the following represents the correct hierarchical order of biological organization from the largest scale to the smallest?}
\begin{enumerate}[label=\Alph*.]
    \item Biosphere $\rightarrow$ Ecosystem $\rightarrow$ Population $\rightarrow$ Community $\rightarrow$ Organism
    \item Biosphere $\rightarrow$ Ecosystem $\rightarrow$ Community $\rightarrow$ Population $\rightarrow$ Organism
    \item Ecosystem $\rightarrow$ Biosphere $\rightarrow$ Community $\rightarrow$ Organism $\rightarrow$ Population
    \item Biosphere $\rightarrow$ Community $\rightarrow$ Ecosystem $\rightarrow$ Population $\rightarrow$ Organism
    \item Ecosystem $\rightarrow$ Community $\rightarrow$ Biosphere $\rightarrow$ Population $\rightarrow$ Organism
\end{enumerate}
\textbf{Answer:} B \\
\textbf{Explanation:} The hierarchy goes from Biosphere to Ecosystems, then Communities (all interacting species in an area), then Populations (individuals of one species), and finally individual Organisms.

\item \textbf{A localized group of organisms that belong to the same species is called a:}
\begin{enumerate}[label=\Alph*.]
    \item Community
    \item Ecosystem
    \item Biosphere
    \item Population
    \item Family
\end{enumerate}
\textbf{Answer:} D \\
\textbf{Explanation:} A population consists of all the individuals of a species living within the bounds of a specified area.

\item \textbf{The idea that novel properties arise at each higher level of biological organization due to the arrangement and interactions of parts is known as:}
\begin{enumerate}[label=\Alph*.]
    \item Evolutionary adaptation
    \item Systems biology
    \item Emergent properties
    \item Reductionism
    \item Feedback regulation
\end{enumerate}
\textbf{Answer:} C \\
\textbf{Explanation:} Emergent properties are absent from the preceding level and appear as complexity increases, like consciousness emerging from the interaction of neurons.

\item \textbf{To understand the chemical basis of inheritance, Watson and Crick studied the molecular structure of DNA. This is an example of:}
\begin{enumerate}[label=\Alph*.]
    \item Systems biology
    \item Reductionism
    \item Deductive reasoning
    \item Emergent properties
    \item Positive feedback
\end{enumerate}
\textbf{Answer:} B \\
\textbf{Explanation:} Reductionism involves breaking down complex biological systems into their simpler, more manageable component parts.

\item \textbf{Which approach to biology aims to model the dynamic behavior of whole biological systems based on the study of the interactions among the system's parts?}
\begin{enumerate}[label=\Alph*.]
    \item Reductionism
    \item Taxonomy
    \item Systems biology
    \item Molecular biology
    \item Bioinformatics
\end{enumerate}
\textbf{Answer:} C \\
\textbf{Explanation:} Systems biology attempts to understand how the components work together to form a functioning whole, complementing reductionism.

\item \textbf{The structural and functional unit of life is the:}
\begin{enumerate}[label=\Alph*.]
    \item DNA molecule
    \item Organ
    \item Atom
    \item Cell
    \item Organism
\end{enumerate}
\textbf{Answer:} D \\
\textbf{Explanation:} The cell is the smallest unit of organization that can perform all activities required for life.

\item \textbf{Which of the following attributes is true of both prokaryotic and eukaryotic cells?}
\begin{enumerate}[label=\Alph*.]
    \item Both contain a nucleus enclosed by a membrane.
    \item Both contain membrane-bound organelles.
    \item Both use DNA to store genetic information.
    \item Both are typically found in multicellular organisms.
    \item Both have an identical cell size.
\end{enumerate}
\textbf{Answer:} C \\
\textbf{Explanation:} All cells, regardless of whether they are prokaryotic or eukaryotic, share certain characteristics like using DNA as their genetic material.

\item \textbf{Which of the following domains contains organisms with prokaryotic cells?}
\begin{enumerate}[label=\Alph*.]
    \item Eukarya and Archaea
    \item Archaea and Bacteria
    \item Bacteria and Protista
    \item Plantae and Fungi
    \item Fungi and Bacteria
\end{enumerate}
\textbf{Answer:} B \\
\textbf{Explanation:} Domain Bacteria and Domain Archaea both consist exclusively of single-celled prokaryotic organisms.

\item \textbf{Genetic information is encoded in a DNA molecule by:}
\begin{enumerate}[label=\Alph*.]
    \item The sequence of its nucleotides (A, T, C, G).
    \item The three-dimensional shape of the double helix.
    \item The exact number of chromosomes in the cell.
    \item The types of proteins that bind to the DNA.
    \item The ratio of sugars to phosphates in the backbone.
\end{enumerate}
\textbf{Answer:} A \\
\textbf{Explanation:} Specific sequences of the four nucleotides (A, T, C, G) act as a chemical alphabet that encodes the information in genes.

\item \textbf{The process by which the information in a gene directs the synthesis of a protein is called:}
\begin{enumerate}[label=\Alph*.]
    \item Replication
    \item Gene expression
    \item Negative feedback
    \item Genomics
    \item Bioinformatics
\end{enumerate}
\textbf{Answer:} B \\
\textbf{Explanation:} Gene expression is the entire process from transcription of DNA to mRNA, to the translation of mRNA into an amino acid chain that forms a protein.

\item \textbf{The universality of the genetic code suggests that:}
\begin{enumerate}[label=\Alph*.]
    \item All life on Earth shares a common evolutionary ancestor.
    \item All organisms have the exact same number of genes.
    \item Prokaryotes and eukaryotes evolved entirely independently.
    \item Evolution by natural selection has ceased.
    \item Every organism expresses all of its genes simultaneously.
\end{enumerate}
\textbf{Answer:} A \\
\textbf{Explanation:} Because essentially the same genetic code is used across all forms of life, it provides strong evidence for common descent.

\item \textbf{Which field of study involves analyzing whole sets of genes and their interactions rather than investigating a single gene at a time?}
\begin{enumerate}[label=\Alph*.]
    \item Proteomics
    \item Genomics
    \item Ecology
    \item Taxonomy
    \item Anatomy
\end{enumerate}
\textbf{Answer:} B \\
\textbf{Explanation:} Genomics is the large-scale analysis of genomes (entire sets of DNA sequences) in one or more species.

\item \textbf{The use of computational tools to store, organize, and analyze the huge volume of biological data is termed:}
\begin{enumerate}[label=\Alph*.]
    \item Proteomics
    \item Systems biology
    \item High-throughput technology
    \item Bioinformatics
    \item Genomics
\end{enumerate}
\textbf{Answer:} D \\
\textbf{Explanation:} Bioinformatics uses computers and mathematics to deal with massive biological datasets resulting from high-throughput sequencing.

\item \textbf{In an ecosystem, energy generally enters as \_\_\_\_\_\_\_ and exits as \_\_\_\_\_\_\_.}
\begin{enumerate}[label=\Alph*.]
    \item heat; light
    \item light; heat
    \item chemical energy; light
    \item light; chemical energy
    \item heat; chemical energy
\end{enumerate}
\textbf{Answer:} B \\
\textbf{Explanation:} Energy flows in a one-way direction through ecosystems, entering mostly as sunlight and exiting as heat lost from metabolic processes.

\item \textbf{Unlike energy, chemical nutrients in an ecosystem:}
\begin{enumerate}[label=\Alph*.]
    \item Are lost to the atmosphere as gases.
    \item Flow in a single, one-way direction.
    \item Are continuously recycled.
    \item Are generated constantly by consumers.
    \item Cannot be utilized by decomposers.
\end{enumerate}
\textbf{Answer:} C \\
\textbf{Explanation:} Chemicals cycle between organisms and the physical environment, continually being used and recycled by producers, consumers, and decomposers.

\item \textbf{When your blood glucose level rises, your pancreas secretes insulin. Insulin stimulates your cells to take up glucose, lowering the blood glucose level, which then shuts off insulin secretion. This is an example of:}
\begin{enumerate}[label=\Alph*.]
    \item Positive feedback
    \item Emergent properties
    \item Negative feedback
    \item Adaptive radiation
    \item Evolutionary adaptation
\end{enumerate}
\textbf{Answer:} C \\
\textbf{Explanation:} In negative feedback, the response (lowering glucose) reduces the initial stimulus (high glucose), turning the pathway off.

\item \textbf{Which of the following is an example of positive feedback in a biological system?}
\begin{enumerate}[label=\Alph*.]
    \item A high concentration of product inhibits the enzyme that produced it.
    \item Sweating cools the body, shutting off the sweating response.
    \item Platelets accumulate at a wound and release chemicals that attract more platelets.
    \item Insulin lowers blood sugar, reducing further insulin release.
    \item A drop in blood pressure triggers a decrease in heart rate.
\end{enumerate}
\textbf{Answer:} C \\
\textbf{Explanation:} In positive feedback, an end product or response speeds up its own production, as seen in the accelerating accumulation of platelets for blood clotting.

\item \textbf{A tree taking in carbon dioxide from the atmosphere and returning oxygen represents which biological theme?}
\begin{enumerate}[label=\Alph*.]
    \item Organisms interact continuously with their environments.
    \item Evolution accounts for life's unity.
    \item The cell is the basic unit of life.
    \item New properties emerge at successive levels.
    \item Life requires transmission of genetic information.
\end{enumerate}
\textbf{Answer:} A \\
\textbf{Explanation:} This gas exchange is a physical interaction between the organism and its environment that changes both.

\item \textbf{Current climate change models predict that human activity, mainly the burning of fossil fuels, will lead to:}
\begin{enumerate}[label=\Alph*.]
    \item A decrease in atmospheric carbon dioxide.
    \item Cooling of the Earth's average temperature.
    \item An increase in global temperatures by at least $3^\circ\text{C}$ before the end of the century.
    \item Increased habitat availability for polar bears.
    \item A reduction in extreme weather events.
\end{enumerate}
\textbf{Answer:} C \\
\textbf{Explanation:} Human-added $CO_2$ acts as a greenhouse gas, trapping heat and raising the average global temperature, causing extensive ecological shifts.

\item \textbf{Biologists have currently identified and named approximately how many living species?}
\begin{enumerate}[label=\Alph*.]
    \item 100,000
    \item 1.8 million
    \item 10 million
    \item 100 million
    \item 1 billion
\end{enumerate}
\textbf{Answer:} B \\
\textbf{Explanation:} About 1.8 million species have been formally identified, though total estimates range from 10 to 100 million.

\item \textbf{Which eukaryotic kingdom consists primarily of multicellular organisms that ingest other organisms?}
\begin{enumerate}[label=\Alph*.]
    \item Plantae
    \item Fungi
    \item Protists
    \item Animalia
    \item Archaea
\end{enumerate}
\textbf{Answer:} D \\
\textbf{Explanation:} Members of the kingdom Animalia obtain food by eating and digesting other organisms.

\item \textbf{Organisms in the kingdom Fungi are distinct because they:}
\begin{enumerate}[label=\Alph*.]
    \item Perform photosynthesis to make their own food.
    \item Are exclusively single-celled prokaryotes.
    \item Absorb nutrients in dissolved form from their surroundings.
    \item Ingest whole organisms and digest them internally.
    \item Live only in extreme environments like hot springs.
\end{enumerate}
\textbf{Answer:} C \\
\textbf{Explanation:} Fungi act mostly as decomposers or symbionts that secrete enzymes to absorb external nutrients.

\item \textbf{Which group of eukaryotes is the most numerous and diverse, mostly consisting of single-celled organisms?}
\begin{enumerate}[label=\Alph*.]
    \item Plants
    \item Animals
    \item Protists
    \item Fungi
    \item Archaea
\end{enumerate}
\textbf{Answer:} C \\
\textbf{Explanation:} Protists are extremely diverse, primarily unicellular eukaryotes that are now classified into several subgroups.

\item \textbf{The shared architecture of cilia in organisms as diverse as Paramecium and human windpipe cells is an example of:}
\begin{enumerate}[label=\Alph*.]
    \item Adaptive radiation
    \item The diversity of life
    \item Unity underlying the diversity of life
    \item Positive feedback
    \item Negative feedback
\end{enumerate}
\textbf{Answer:} C \\
\textbf{Explanation:} Shared structural features like the internal tubule arrangement in cilia indicate a common evolutionary origin despite wide species differences.

\item \textbf{The scientific explanation for the unity and diversity of organisms is:}
\begin{enumerate}[label=\Alph*.]
    \item The flow of energy
    \item Genomics
    \item Reductionism
    \item Evolution
    \item Gene expression
\end{enumerate}
\textbf{Answer:} D \\
\textbf{Explanation:} Evolution is the core theme of biology, explaining how species accumulate differences from ancestors while retaining shared, foundational traits.

\item \textbf{Charles Darwin's groundbreaking book published in 1859 is titled:}
\begin{enumerate}[label=\Alph*.]
    \item On the Origin of Species by Means of Natural Selection
    \item The Descent of Man
    \item Principles of Geology
    \item The Double Helix
    \item Systems Biology and Ecology
\end{enumerate}
\textbf{Answer:} A \\
\textbf{Explanation:} Darwin's famous book articulated his concepts of descent with modification and natural selection.

\item \textbf{Darwin's phrase "descent with modification" accurately captures:}
\begin{enumerate}[label=\Alph*.]
    \item The flow of energy from sun to plants.
    \item The duality of life's unity and diversity.
    \item The mechanism of gene translation.
    \item The cycling of nutrients in an ecosystem.
    \item The creation of entirely new domains of life.
\end{enumerate}
\textbf{Answer:} B \\
\textbf{Explanation:} Unity comes from kinship of descent from common ancestors, and diversity from the modifications evolved as they branched out.

\item \textbf{Which of the following was NOT one of Darwin's three key observations from nature?}
\begin{enumerate}[label=\Alph*.]
    \item Individuals in a population vary in their traits.
    \item Traits are generally heritable.
    \item A population can produce far more offspring than can survive.
    \item DNA is the genetic material responsible for heredity.
    \item Species generally are suited to their environments.
\end{enumerate}
\textbf{Answer:} D \\
\textbf{Explanation:} Darwin did not know about DNA; his theory was based on phenotypic observations, overpopulation, and apparent environmental fit.

\item \textbf{Natural selection occurs because:}
\begin{enumerate}[label=\Alph*.]
    \item Organisms consciously adapt to their environments.
    \item The environment creates new mutations on demand.
    \item Individuals with traits better suited to the environment survive and reproduce more successfully.
    \item Populations produce exactly the number of offspring the environment can support.
    \item Evolution always progresses toward perfection.
\end{enumerate}
\textbf{Answer:} C \\
\textbf{Explanation:} The unequal reproductive success of individuals with varying heritable traits is the driving force of natural selection.

\item \textbf{The skeletal differences in the forelimbs of a bat, a human, and a whale represent:}
\begin{enumerate}[label=\Alph*.]
    \item Convergent evolution from different ancestors.
    \item The exact same sequence of DNA.
    \item The independent emergence of bone structures.
    \item Anatomical variations inherited from a common mammalian ancestor.
    \item Structural defects that decrease fitness.
\end{enumerate}
\textbf{Answer:} D \\
\textbf{Explanation:} These structures are homologous; they represent modifications of a common architecture adapted for different environments.

\item \textbf{The radiation of the Galápagos finches from a single ancestral species to multiple distinct species with different beak shapes is an example of:}
\begin{enumerate}[label=\Alph*.]
    \item A single gene mutation
    \item A lack of genetic variation
    \item Descent with modification and adaptation to different environments
    \item Energy cycling
    \item Negative feedback
\end{enumerate}
\textbf{Answer:} C \\
\textbf{Explanation:} The finches branched into various species as isolated populations adapted to different food sources over many generations.

\item \textbf{In scientific inquiry, recorded observations are specifically called:}
\begin{enumerate}[label=\Alph*.]
    \item Hypotheses
    \item Theories
    \item Conclusions
    \item Data
    \item Variables
\end{enumerate}
\textbf{Answer:} D \\
\textbf{Explanation:} Data are the raw items of information, either qualitative or quantitative, on which scientific inquiry is based.

\item \textbf{Jane Goodall's decades of recording the behavior of chimpanzees primarily produced which type of data initially?}
\begin{enumerate}[label=\Alph*.]
    \item Quantitative
    \item Controlled
    \item Experimental
    \item Qualitative
    \item Computational
\end{enumerate}
\textbf{Answer:} D \\
\textbf{Explanation:} Qualitative data consist of recorded descriptions, such as Goodall's field notebooks of chimpanzee behavior.

\item \textbf{Deriving a general principle from a large number of specific observations is a process called:}
\begin{enumerate}[label=\Alph*.]
    \item Deductive reasoning
    \item Inductive reasoning
    \item Experimental control
    \item Hypothesis testing
    \item Statistical modeling
\end{enumerate}
\textbf{Answer:} B \\
\textbf{Explanation:} Inductive reasoning flows from specific observations to broad generalizations (e.g., "All living things are made of cells").

\item \textbf{A scientific hypothesis must be:}
\begin{enumerate}[label=\Alph*.]
    \item Proven true beyond all doubt.
    \item Broad enough to explain all phenomena.
    \item Accepted by societal consensus before testing.
    \item Testable and falsifiable.
    \item Derived using deductive reasoning only.
\end{enumerate}
\textbf{Answer:} D \\
\textbf{Explanation:} A valid hypothesis leads to predictions that can be tested observationally or experimentally, with the potential to be proven false.

\item \textbf{When researchers use "If... then" logic to predict specific results from a general premise, they are using:}
\begin{enumerate}[label=\Alph*.]
    \item Inductive reasoning
    \item Deductive reasoning
    \item Emergent properties
    \item Systems biology
    \item Bioinformatics
\end{enumerate}
\textbf{Answer:} B \\
\textbf{Explanation:} Deductive reasoning extrapolates specific expected outcomes from general premises if a hypothesis is true.

\item \textbf{In an experiment testing the effect of camouflage on mouse predation rates, what was the independent variable?}
\begin{enumerate}[label=\Alph*.]
    \item The amount of moonlight.
    \item The rate of predation.
    \item The color pattern of the mouse models.
    \item The type of predator (bird vs. mammal).
    \item The type of soil in the habitat.
\end{enumerate}
\textbf{Answer:} C \\
\textbf{Explanation:} The independent variable is the factor manipulated by the researchers, which in this case was the coloration of the mouse models.

\item \textbf{In the same mouse camouflage experiment, what was the dependent variable?}
\begin{enumerate}[label=\Alph*.]
    \item The color of the mouse models.
    \item The number of models placed in the environment.
    \item The amount of predation (attack rate) on the models.
    \item The exact geographical location of the test.
    \item The temperature during the night.
\end{enumerate}
\textbf{Answer:} C \\
\textbf{Explanation:} The dependent variable is the factor being measured, which is predicted to be affected by the independent variable.

\item \textbf{What is the primary purpose of a control group in a scientific experiment?}
\begin{enumerate}[label=\Alph*.]
    \item To ensure that all variables in the environment are physically locked down.
    \item To provide a standard of comparison to cancel out the effects of unwanted variables.
    \item To test a secondary hypothesis simultaneously.
    \item To prove the hypothesis is true.
    \item To introduce a new independent variable.
\end{enumerate}
\textbf{Answer:} B \\
\textbf{Explanation:} A control group serves as a baseline that differs from the experimental group by only the independent variable, isolating its effect.

\item \textbf{Can a scientific hypothesis be proven absolutely true?}
\begin{enumerate}[label=\Alph*.]
    \item Yes, if tested mathematically.
    \item Yes, if it is tested more than ten times.
    \item No, because it is impossible to test all possible alternative hypotheses.
    \item No, because data in biology are entirely subjective.
    \item Yes, once it is published in a peer-reviewed journal.
\end{enumerate}
\textbf{Answer:} C \\
\textbf{Explanation:} Testing supports a hypothesis by failing to prove it incorrect, but one can never definitively eliminate every possible alternative explanation.

\item \textbf{Why does science inherently ignore hypotheses involving invisible ghosts or supernatural spirits?}
\begin{enumerate}[label=\Alph*.]
    \item They are too mathematically complex.
    \item Science only deals with natural, testable explanations for natural phenomena.
    \item Scientists have already proven spirits do not exist.
    \item It is unethical to experiment on supernatural entities.
    \item Such hypotheses are strictly inductive rather than deductive.
\end{enumerate}
\textbf{Answer:} B \\
\textbf{Explanation:} Science requires testability. Supernatural entities cannot be tested with empirical data, placing them outside the bounds of science.

\item \textbf{How does a scientific theory differ from a scientific hypothesis?}
\begin{enumerate}[label=\Alph*.]
    \item A theory is a specific, narrow guess; a hypothesis is a broad fact.
    \item A theory is broader in scope, can generate new hypotheses, and is supported by a massive body of evidence.
    \item A theory can never be rejected, whereas a hypothesis can.
    \item A hypothesis is an unproven theory.
    \item Theories are strictly qualitative; hypotheses are quantitative.
\end{enumerate}
\textbf{Answer:} B \\
\textbf{Explanation:} Theories (like the Theory of Natural Selection) are comprehensive, evidence-backed frameworks that encompass many hypotheses.

\item \textbf{Which of the following best states the basic difference between science and technology?}
\begin{enumerate}[label=\Alph*.]
    \item Science focuses on applied research; technology focuses on pure research.
    \item Science aims to understand natural phenomena; technology aims to apply scientific knowledge for a specific purpose.
    \item Science is performed only by individuals; technology is a team effort.
    \item Science relies on qualitative data; technology relies on quantitative data.
    \item Science deals with facts; technology deals with theories.
\end{enumerate}
\textbf{Answer:} B \\
\textbf{Explanation:} Science discovers how the world works, while technology uses those discoveries to invent tools and solve societal problems.

\item \textbf{An organism that is easy to grow in the lab and lends itself well to investigations whose findings can be broadly applied is known as a:}
\begin{enumerate}[label=\Alph*.]
    \item Model organism
    \item Prototype species
    \item Control group
    \item Emergent organism
    \item Eukaryotic standard
\end{enumerate}
\textbf{Answer:} A \\
\textbf{Explanation:} Model organisms like \textit{Drosophila} (fruit fly) or \textit{E. coli} share basic biological processes with many species, making them ideal for laboratory research.

\item \textbf{Why is diversity in the scientific community considered important?}
\begin{enumerate}[label=\Alph*.]
    \item It is required by international law.
    \item It guarantees that all hypotheses are eventually proven true.
    \item Diverse viewpoints lead to more robust, valuable, and productive scientific interchange and innovation.
    \item It eliminates the need for the peer-review process.
    \item It prevents any two scientists from working on the same problem.
\end{enumerate}
\textbf{Answer:} C \\
\textbf{Explanation:} Bringing different backgrounds and perspectives to scientific problems fosters creativity and helps uncover varied technological solutions.

\item \textbf{If a researcher evaluates the DNA sequences of an entire species' genome, this field of study is referred to as:}
\begin{enumerate}[label=\Alph*.]
    \item Physiology
    \item Genomics
    \item Systems Biology
    \item Proteomics
    \item Anatomy
\end{enumerate}
\textbf{Answer:} B \\
\textbf{Explanation:} Genomics is the study of whole sets of genes and their interactions within an organism's genome.

\item \textbf{The flow of genetic information within a cell generally proceeds as:}
\begin{enumerate}[label=\Alph*.]
    \item RNA $\rightarrow$ DNA $\rightarrow$ Protein
    \item Protein $\rightarrow$ RNA $\rightarrow$ DNA
    \item DNA $\rightarrow$ RNA $\rightarrow$ Protein
    \item DNA $\rightarrow$ Protein $\rightarrow$ RNA
    \item RNA $\rightarrow$ Protein $\rightarrow$ DNA
\end{enumerate}
\textbf{Answer:} C \\
\textbf{Explanation:} During gene expression, DNA is transcribed into mRNA, which is then translated into an amino acid sequence (protein).

\item \textbf{Which statement about feedback regulation is accurate?}
\begin{enumerate}[label=\Alph*.]
    \item Positive feedback loops are more common in biological systems than negative feedback loops.
    \item In negative feedback, the response amplifies the initial stimulus.
    \item Biological systems exclusively rely on positive feedback.
    \item In negative feedback, the end product inhibits the pathway that produces it.
    \item Feedback regulation only occurs at the ecosystem level.
\end{enumerate}
\textbf{Answer:} D \\
\textbf{Explanation:} Negative feedback is the primary mechanism of biological regulation; the accumulation of an end product acts to slow its own production.

\item \textbf{A team of researchers publishes their experimental results. Before publication, their paper is evaluated by independent, anonymous experts in the same field. This process is called:}
\begin{enumerate}[label=\Alph*.]
    \item Deductive reasoning
    \item Experimental control
    \item Bioinformatics
    \item Peer review
    \item Emergent feedback
\end{enumerate}
\textbf{Answer:} D \\
\textbf{Explanation:} Peer review ensures the integrity and validity of scientific claims before they become part of the scientific literature.

\item \textbf{In the context of evolutionary theory, what does "adaptation" refer to?}
\begin{enumerate}[label=\Alph*.]
    \item The conscious effort an organism makes to survive its environment.
    \item The immediate physical changes an organism undergoes during its lifespan.
    \item Inherited characteristics of organisms that enhance their survival and reproduction in specific environments.
    \item The survival of an organism regardless of its traits.
    \item The artificial selection of traits by humans.
\end{enumerate}
\textbf{Answer:} C \\
\textbf{Explanation:} Evolutionary adaptation refers to heritable traits that have accumulated over generations because they increase fitness in a particular environment.

\end{enumerate}

\end{document}
